\section{Orientation-Relationship Index Correspondence}

An orientation relationship (OR) between a parent phase $p$ and a child phase $c$ stores the
rigid rotation $\mathbf{R}$ that carries parent crystal-frame Cartesian vectors into the child
crystal frame. Rotations act on Cartesian vectors, not on Miller indices: pushing $[uvw]$ or
$(hkl)$ triples through $\mathbf{R}$ directly is only valid for cubic phases with equal lattice
parameters, and silently wrong otherwise. This note fixes the index-correspondence algorithm
implemented in \texttt{pytex.core.transformation}.

\subsection{Correspondence Matrices}

Let $\mathbf{A}_{p}$ and $\mathbf{A}_{c}$ be the direct structure matrices of the two phases
(columns are the lattice basis vectors expressed in the respective crystal Cartesian frames), and
let $\mathbf{B} = \mathbf{A}^{-\mathsf{T}}$ be the corresponding reciprocal structure matrices
under the PyTex normalization $\mathbf{a}^{*}_{i} \cdot \mathbf{a}_{j} = \delta_{ij}$.

A direction with components $\mathbf{u}_{p}$ (a $[uvw]$ triple) has the Cartesian image
$\mathbf{A}_{p}\mathbf{u}_{p}$; its image in the child crystal frame is
$\mathbf{R}\,\mathbf{A}_{p}\mathbf{u}_{p}$, and reading that vector back in child lattice
coordinates gives the \emph{direction-index correspondence}

\begin{equation}
\mathbf{u}_{c} = \mathbf{M}\,\mathbf{u}_{p},
\qquad
\mathbf{M} = \mathbf{A}_{c}^{-1}\,\mathbf{R}\,\mathbf{A}_{p}.
\end{equation}

Plane indices transform through the reciprocal bases:

\begin{equation}
\mathbf{h}_{c} = \mathbf{M}^{*}\,\mathbf{h}_{p},
\qquad
\mathbf{M}^{*} = \mathbf{B}_{c}^{-1}\,\mathbf{R}\,\mathbf{B}_{p}
= \mathbf{A}_{c}^{\mathsf{T}}\,\mathbf{R}\,\mathbf{A}_{p}^{-\mathsf{T}}
= \mathbf{M}^{-\mathsf{T}},
\end{equation}

where the last equality uses $\mathbf{R}^{-\mathsf{T}} = \mathbf{R}$ for a proper rotation. The
inverse-transpose relation implies zone-law invariance,

\begin{equation}
\mathbf{h}_{c} \cdot \mathbf{u}_{c}
= \mathbf{h}_{p}^{\mathsf{T}} \mathbf{M}^{-1} \mathbf{M} \, \mathbf{u}_{p}
= \mathbf{h}_{p} \cdot \mathbf{u}_{p},
\end{equation}

so a direction lying in a plane maps to a direction lying in the mapped plane. Both identities
are pinned by unit tests. Per-variant correspondences replace $\mathbf{R}$ by the variant
rotation $\mathbf{R}\,\mathbf{S}_{p}^{\mathsf{T}}$.

\subsection{Rationalization}

$\mathbf{M}$ and $\mathbf{M}^{*}$ are generally irrational, so the exact image of an integer
triple is generally not an integer triple. PyTex reports both: the exact components, and the
nearest primitive integer triple $\mathbf{n}$ within a bound $\max_i |n_i| \le N$
(default $N = 17$, covering the Greninger--Troiano $\langle 5\,12\,17 \rangle$ family). The
comparison metric is angular and sign-sensitive, evaluated in the correct space: candidate
triples are compared through their Cartesian images ($\mathbf{A}\mathbf{n}$ for directions,
$\mathbf{B}\mathbf{n}$ for plane normals), and the winner minimizes the angle to the exact
image. The residual is computed with the $\operatorname{atan2}$ form,

\begin{equation}
\delta = \operatorname{atan2}\!\left(
  \lVert \hat{\mathbf{x}} \times \hat{\mathbf{t}} \rVert,\;
  \hat{\mathbf{x}} \cdot \hat{\mathbf{t}}
\right),
\end{equation}

which keeps full floating-point precision near $\delta = 0$ where $\arccos$ saturates. A zero
residual therefore certifies an exact crystallographic parallelism, and a nonzero residual
quantifies "nearly parallel" instead of hiding it in a rounding step.

The candidate set enumerates all signed primitive triples ($\gcd = 1$) with entries in
$[-N, N]$; the set is cached per bound and the search is a single vectorized matrix product, so
the cost is negligible against the surrounding workflow.

\subsection{Failure Modes And Limits}

\begin{itemize}
  \item The rationalization bound truncates the search: an OR whose true correspondence involves
        indices beyond $N$ reports the best in-bound triple with a correspondingly large
        residual. The bound is a keyword argument on every mapping call.
  \item Correspondence matrices assume both phases' structure matrices are expressed in their
        canonical crystal Cartesian frames; construction-time phase checks enforce membership.
  \item The deformation gradient (Bain strain) is deliberately \emph{not} part of this surface;
        index correspondence is purely rotational plus metric, and strain analysis will land as
        a separate, explicitly named feature.
\end{itemize}

\subsection*{Normative references}

International Tables for Crystallography, Vol.~A (basis and reciprocal-basis conventions).

\subsection*{Informative references}

Kurdjumov, G., Sachs, G., Z.~Phys.~64 (1930) 325.
Bain, E.~C., Trans.~AIME 70 (1924) 25.
Bhadeshia, H.~K.~D.~H., \emph{Geometry of Crystals, Polycrystals, and Phase Transformations}.
